\documentclass[11pt,a4paper]{article}

%  Idioma, codificación y página
\usepackage[spanish]{babel}
\usepackage[T1]{fontenc}
\usepackage[a4paper,margin=3.2cm,headheight=14pt,headsep=1cm]{geometry}

%  Matemática y utilidades
\usepackage{amssymb,amsmath,amsthm,amsfonts}
\usepackage{mathrsfs}
\usepackage{mathtools}
\usepackage{xcolor}
\usepackage{graphicx}
\usepackage{booktabs}
\usepackage{tikz-cd}
\usepackage[shortlabels,inline]{enumitem}

%  Tipografía: texto y matemática en Times (como Lyons--Peres)
\IfFileExists{newtxtext.sty}{%
  \usepackage{newtxtext}%
  \usepackage{newtxmath}%
}{%
  \usepackage{mathptmx}%
}
\usepackage{esint}
\usepackage{microtype}

% Colores: gris para los números, verde oscuro para los enlaces
\definecolor{LPgray}{gray}{0.55}
\definecolor{LPgreen}{RGB}{6,92,10}

% Fuente de los títulos: Times negrita, condensada horizontalmente
\newcommand{\LPscale}{0.8}
\newlength{\LPw}
% ancho para las cajas condensadas = (5/4)(argumento), solo con enteros
\newcommand{\LPsetwidth}[1]{\setlength{\LPw}{\dimexpr 5\dimexpr #1\relax/4\relax}}
\newcommand{\LPhead}{\normalfont\bfseries}

%  Atajos: conjuntos 
\newcommand{\N}{\mathbb{N}}
\newcommand{\Z}{\mathbb{Z}}
\newcommand{\Q}{\mathbb{Q}}
\newcommand{\R}{\mathbb{R}}
\newcommand{\C}{\mathbb{C}}
\newcommand{\F}{\mathbb{F}}
\newcommand{\eps}{\varepsilon}
\newcommand{\ds}{\displaystyle}

%  Atajos: proba
\newcommand{\Prob}{\textbf{P}}          
\newcommand{\E}{\textbf{E}}             
\newcommand{\ind}[1]{\mathbf{1}_{#1}}   
\DeclareMathOperator{\Var}{Var}
\DeclareMathOperator{\Cov}{Cov}
\DeclareMathOperator{\Corr}{Corr}
\newcommand{\dto}{\xrightarrow{\;d\;}}        % convergencia en distribución
\newcommand{\pto}{\xrightarrow{\;\Prob\;}}    % convergencia en probabilidad
\newcommand{\asto}{\xrightarrow{\;\mathrm{c.s.}\;}}  % casi segura
\newcommand{\iid}{\text{i.i.d.}}

%  Más
\DeclareMathOperator{\supp}{supp}
\DeclareMathOperator{\diag}{diag}
\DeclareMathOperator{\tr}{tr}
\DeclareMathOperator{\rank}{rank}
\DeclareMathOperator{\im}{im}
\DeclareMathOperator{\sgn}{sgn}
\DeclareMathOperator{\id}{id}
\DeclareMathOperator*{\argmin}{arg\,min}
\DeclareMathOperator*{\argmax}{arg\,max}
\newcommand{\dd}{\,\mathrm{d}}                          % \int f \dd x
\newcommand{\deriv}[2]{\frac{\mathrm{d}#1}{\mathrm{d}#2}}
\newcommand{\pderiv}[2]{\frac{\partial #1}{\partial #2}}
\newcommand{\defeq}{\coloneqq}                           % :=
\newcommand{\inv}{^{-1}}
\newcommand{\cl}[1]{\overline{#1}}                       % clausura

% Delimitadores que se ajustan solos (con asterisco crecen: \abs*{...})
\DeclarePairedDelimiter{\abs}{\lvert}{\rvert}
\DeclarePairedDelimiter{\norm}{\lVert}{\rVert}
\DeclarePairedDelimiter{\floor}{\lfloor}{\rfloor}
\DeclarePairedDelimiter{\ceil}{\lceil}{\rceil}
\DeclarePairedDelimiter{\set}{\lbrace}{\rbrace}
\newcommand{\inner}[2]{\left\langle #1,\, #2\right\rangle}

% Nota para uno mismo, bien visible
\newcommand{\todo}[1]{{\color{red}[\textbf{TODO:} #1]}}

% Término que se define: negrita cursiva, como en el libro
\newcommand{\defn}[1]{\textbf{\emph{#1}}}

%  Secciones: centradas, con el número en gris
\usepackage{titlesec}
\titleformat{\section}[block]{\centering\large\bfseries}%
  {\textcolor{LPgray}{\thesection}\hspace{0.6em}}{0pt}{}
\titleformat{name=\section,numberless}[block]{\centering\large\bfseries}{}{0pt}{}
\titlespacing*{\section}{0pt}{3.2ex plus 1ex minus .5ex}{1.8ex plus .3ex}

\titleformat{\subsection}[block]{\bfseries}%
  {\textcolor{LPgray}{\thesubsection}\hspace{0.6em}}{0pt}{}
\titlespacing*{\subsection}{0pt}{2.6ex plus 1ex minus .5ex}{1.2ex plus .2ex}

\titleformat{\paragraph}[runin]{\bfseries\itshape}{}{0pt}{}[.]
\titlespacing*{\paragraph}{0pt}{1.5ex plus .5ex}{0.6em}

%  Encabezados y pies de página
\usepackage{fancyhdr}
\renewcommand{\sectionmark}[1]{\markright{\thesection\hspace{0.6em}#1}}
\pagestyle{fancy}
\fancyhf{}
\renewcommand{\headrulewidth}{0pt}
\fancyhead[L]{\small\textit{\rightmark}}
\fancyhead[R]{\small\textbf{\thepage}}
\fancypagestyle{plain}{%
  \fancyhf{}%
  \renewcommand{\headrulewidth}{0pt}%
  \fancyhead[R]{\small\textbf{\thepage}}%
}

%  Enlaces en verde y referencias cruzadas
\usepackage{hyperref}
\hypersetup{colorlinks=true,linkcolor=LPgreen,citecolor=LPgreen,
  urlcolor=LPgreen}
% cleveref: \cref{etiqueta} escribe "Teorema 3.1", "Definición 2.1", etc.,
% aunque los entornos compartan contador. Plural automático con varias etiquetas.
\usepackage[nameinlink]{cleveref}
\crefname{teo}{Teorema}{Teoremas}
\crefname{lem}{Lema}{Lemas}
\crefname{prop}{Proposición}{Proposiciones}
\crefname{coro}{Corolario}{Corolarios}
\crefname{conj}{Conjetura}{Conjeturas}
\crefname{defi}{Definición}{Definiciones}
\crefname{obs}{Observación}{Observaciones}
\crefname{ejemplo}{Ejemplo}{Ejemplos}
\crefname{notacion}{Notación}{Notaciones}
\crefname{ej}{Ejercicio}{Ejercicios}
\crefname{section}{Sección}{Secciones}
\crefname{subsection}{Subsección}{Subsecciones}
\crefname{equation}{ecuación}{ecuaciones}
\crefname{figure}{Figura}{Figuras}
\crefname{table}{Tabla}{Tablas}
\Crefname{equation}{Ecuación}{Ecuaciones}

%  Teoremas: nombre en negrita, enunciado en cursiva, fin de prueba ◀
\newtheoremstyle{LPplain}%
  {9pt}{9pt}{\itshape}{0pt}{\bfseries}{.}{0.5em}%
  {\thmname{#1}\thmnumber{ #2}\thmnote{ (#3)}}
\newtheoremstyle{LPdef}%
  {9pt}{9pt}{\normalfont}{0pt}{\bfseries}{.}{0.5em}%
  {\thmname{#1}\thmnumber{ #2}\thmnote{ (#3)}}
\newtheoremstyle{LPrem}%
  {9pt}{9pt}{\normalfont}{0pt}{\itshape}{.}{0.5em}%
  {\thmname{#1}\thmnumber{ #2}\thmnote{ (#3)}}

\numberwithin{equation}{section}

\theoremstyle{LPplain}
\newtheorem{teo}{Teorema}[section]
\newtheorem{lem}[teo]{Lema}
\newtheorem{prop}[teo]{Proposición}
\newtheorem{coro}[teo]{Corolario}
\newtheorem{conj}[teo]{Conjetura}
\newtheorem*{teo*}{Teorema}

\theoremstyle{LPdef}
\newtheorem{defi}[teo]{Definición}
\newtheorem{ej}{Ejercicio}[section]
\newtheorem*{ej*}{Ejercicio}

\theoremstyle{LPrem}
\newtheorem{obs}[teo]{Observación}
\newtheorem{ejemplo}[teo]{Ejemplo}
\newtheorem{notacion}[teo]{Notación}

\renewcommand{\qedsymbol}{$\blacktriangleleft$}
\renewcommand{\proofname}{Demostración}

%  Título y resumen
\makeatletter
\renewcommand{\maketitle}{%
  \thispagestyle{plain}%
  \begin{center}
    \textcolor{LPgray}{\rule{\textwidth}{0.4pt}}\par
    \vspace{1.2em}
    \LPsetwidth{\textwidth}%
    \scalebox{\LPscale}[1]{\parbox[t]{\LPw}{\centering
      \LPhead\fontsize{28}{31}\selectfont \@title\par}}\par
    \vspace{1em}
    {\large \@author\par}
    \vspace{0.3em}
    {\normalsize \@date\par}
    \vspace{1em}
    \textcolor{LPgray}{\rule{\textwidth}{0.4pt}}\par
  \end{center}
  \vspace{1em}}

\renewenvironment{abstract}{%
  \begin{center}\begin{minipage}{0.86\textwidth}
  \small\noindent\textbf{\abstractname.}\ }{%
  \end{minipage}\end{center}\vspace{0.5em}}
\makeatother

%  Datos del documento
\title{Paseos al azar y redes eléctricas}
\author{Lolistone
}
\date{\today}

%  Documento

\begin{document}

\maketitle

\begin{abstract}
    Veremos que puede pensarse a las cadenas de Markov reversibles como
    redes eléctricas. Probaremos resultados clásicos de la teoría de Markov
    apoyandonos en la intuición física ganada.
\end{abstract}

% \tableofcontents   % descomentar si el texto es largo

\section{¿Redes eléctricas, cadenas reversibles?}

  Informalmente, una red eléctrica es un circuito de toda la vida, con
  sus leyes y propiedades. Mientras que una cadena reversible es una que 
  podemos mirar tanto de adelante hacia atrás como de atrás hacia adelante,
  de ahí el \emph{reversible}. Sin haberlos definido formalmente, a simple 
  vista no se ve ninguna conexión muy clara, sin embargo existe y nos será
  muy útil para probar resultados acerca de estas cadenas. 

\section{Preliminares}

\begin{defi}[Cadena reversible]
  Una cadena \defn{reversible} es una cadena de Markov que cumple 
  \defn{balance detallado}, es decir, existe $\pi$ una medida tal que
  \begin{equation}\label{eq:bd}
   \pi(x)\Prob(x,y) = \pi(y)\Prob(y,x), \quad \forall x,y \in \mathcal{X}.
  \end{equation}
   donde $\Prob$ es la matriz de transición de la cadena y $\mathcal{X}$ el espacio
  de estados.
\end{defi}

\begin{obs}
  ¿Donde está lo de reversible? Bueno, notemos que la condición de balance detallado
  nos dice que si $X_0 \sim \pi$,
  $$\Prob(X_0 = x, X_1 = y) = \Prob(X_0 = y, X_1 = x), \quad \forall x,y \in \mathcal{X}.$$
  Más aún, via inducción uno puede probar que 
  $$(X_0, X_1, \dots, X_k) \sim (X_k, \cdots, X_1, X_0), \quad \forall k \in \N.$$
  Esto nos dice que ver la película al derecho o al revés es lo mismo.
\end{obs}

\begin{obs}
  \textbf{Ojo}, en ~\ref{eq:bd} $\pi$ puede no ser una distribución. 
\end{obs}

\begin{defi}[Red]\label{def:red}
  Una \defn{red} es un grafo $G=(V,E)$ no orientado junto con
  $ \{ C(e) \}_{e\in E}$ con $c \geq 0$ llamada \defn{conductancia}. Llamamos
  \defn{resistencia} y notamos $\{r(e)\}_{e\in E}$ a su inversa, $r(e) = \frac{1}{c(e)}$.
\end{defi}

\begin{notacion}
  Escribimos $(G, \{C(e)\}_{e\in E})$ para una red.
\end{notacion}

\begin{obs}
  Si bien sobre $G$ no imponemos ninguna condición, los primeros resultados
  que probemos sobre redes van asumir que este es finito.
\end{obs}

Ahora que ya sabemos que es cada cosa veamos que dada una cadena podemos 
obtener una red y viceversa. Comencemos tomando una red $(G, \{C(e)\}_{e\in E})$
y consideremos la matriz de transición $\Prob$ dada por 
$$ \Prob(x, y) = \frac{c(x,y)}{c(x)}, \quad c(x) = \sum_{y \colon y \sim x} c(x, y)$$
donde $y \sim x$ si $y$ es vecino de $x$. A este proceso también lo vamos a llamar
\defn{paseo al azar pesado}. Consideremos ahora
$$ \pi(x) = \frac{c(x)}{c_G}, \quad c_G = \sum_{x \in V} c(x).$$
Es fácil chequear que esta condición cumple \emph{balance detallado} y por lo tanto
tenemos una cadena reversible. Supongamos ahora que tenemos un cadena reversible en un
espacio de estados finito, veamos que podemos representarla como un paseo pesado sobre 
una red, definimos $G$ un grafo con vertices en $\mathcal{X}$, $x \sim y$ si $P(x, y) > 0$
y 
$$c(x, y) = \pi(x)\Prob(x, y), \quad \forall x,y \in \mathcal{X}.$$
Con esta elección de pesos tenemos que $c(x) = \pi(x)$ y la matriz asociada es $P$.

\begin{obs}
  $\pi$ no necesariamente es distribución. Por ejemplo, en un paseo al azar simple podemos
  tomar $\pi(x) = deg(x)$. En este caso, $c(x, y) = 1.$
\end{obs}

\section{Funciones armónicas y tiempos de llegada}

Dados $A, Z$ subconjunto disjuntos del espacio de estados, es de nuestro interés por ejemplo
ver la probabilidad llegar primero a $A$ que a $Z$ dado que comenzamos en un estado cualquiera $x$.
Es decir, nos interesa estudiar
$$F(x) = \Prob_x[\tau_A < \tau_Z].$$
donde $\tau_A = \min\{t \geq 0 \colon X_t \in A\}$. También vamos a considerar 
$\tau_{A}^{+} = \min\{t > 0 \colon X_t \in A\}$ que difiere con el primer tiempo de parada cuando
el proceso comienza en $A$. Veamos que tenemos,
\begin{align*}
  F(x) &= \sum_{y \sim x} \Prob_x[\text{el primer paso es } y]\,
          \Prob_x[\tau_A < \tau_Z \mid \text{el primer paso es } y] \\
       &= \sum_{y \sim x} P(x,y)\, \Prob_y[\tau_A < \tau_Z] \\
       &= \sum_{y \sim x} P(x,y)\, F(y).
\end{align*}
Si la cadena es reversible esto es 
$$ \frac{1}{\pi(x)} \sum_{y \sim x} c(x, y) F(y),$$
y si consideramos un paseo al azar simple obtenemos
$$ \frac{1}{d(x)} \sum_{y \sim x} F(y).$$
Esta última expresión parece una versión discreta de nuestra noción 
de función armónica de toda la vida, esto motiva la siguiente definición.

\begin{defi}[Armónica]
Una función $h \colon \mathcal{X} \to \R$ es \defn{armónica} en $x$ si
$$  h(x) = \sum_{y} P(x,y)\,h(y). $$
\end{defi}

Al igual que con la noción usual vale un \defn{principio del máximo}.

\begin{prop}[Principio del máximo]
  Sea $W$ un conjunto de estados contable de un espacio de estados $\mathcal{X}$.
  Si $h : \mathcal{X} \to \R$ es armónica y su supremo sobre $\mathcal{X}$ se alcanza en algún
  $x_0 \in W$, entonces es constante en todos los estados accesibles desde $W$.
\end{prop}

\begin{proof}
  Sea $K \defeq \{x \in \mathcal{X} \colon h(x) = \sup_{\mathcal{X}}{h} \}$. 
  Tomemos $x \in K \cap W$ si $P(x, y) > 0$ tenemos que $y \in K$ pues si $h(y) < h(x)$ 
  \begin{align*}
    h(x) &= \sum_{z \in \mathcal{X}} P(x, z) h(z) \\ 
         &= h(y) P(x,y) + \sum_{z \in \mathcal{X} \setminus \{y\}} P(x, z) h(z) \\
         &< h(x) P(x,y) + \sum_{z \in \mathcal{X} \setminus \{y\}} P(x, z) h(x) = h(x).
  \end{align*}
  Lo cual es absurdo. Notemos que en el último paso usamos $h(z) \leq h(x)$ y la suposición. 
\end{proof}

\begin{prop}[Principio de unicidad]
  Sea $W$ un subconjunto finito de estado de $\mathcal{X}$ un espacio contable. Supongamos
  que se puede llegar a todo elemento de $\mathcal{X}$ desde $W$. Si $h, g : \mathcal{X} \to \R$
  son armónicas en $W$ y coinciden entonces son iguales.
\end{prop}

\begin{proof} Consideremos $f = h - g$. Cómo $W$ es finito y fuera de este vale $0$
$f$ alcanza máximo en $\mathcal{X}$, si $x \notin W$ entonces $f \leq 0$, y si $x \in W$
por el principio del máximo $f = 0$ en todos lados. Luego, $f \leq 0$. Por simetría, 
mirando $-f$ tenemos que $f \geq 0$ y por lo tanto es identicamente nula. 
\end{proof}

Ahora probemos un resultado que nos va a ser muy útil a futuro.

\begin{prop}[Principio de existencia]\label{prop:extension}
  Sea $(X_t)_t$ una cadena de Markov irreducible, $W \subset \mathcal{X}$ con $\mathcal{X}$ contable 
  y $\mathcal{X} \setminus W$ finito y $h_{W} \colon W \to \R$. La función $h(x) = \E_x h_W(X_{\tau_W})$ es la única
  extensión armónica de $h_W$ a $\mathcal{X}\setminus W$.  
\end{prop}

\begin{proof}
  Por el principio del máximo y el de unicidad basta ver que $h$ es armónica. Es fácil chequear 
  que coinciden en $B$ veamos que ocurre en el complemento. Sea $x \notin B$,
  \begin{align*}
    \E_x h_B(X_{\tau_B}) &= \sum_{y \in X} P(x, y) \E_x[h_B(X_{\tau_B}) \mid X_1 = y] \\
                         &= \sum_{y \in X} P(x, y) \E_y h_B(X_{\tau_B}). 
  \end{align*}
  Luego, $h$ es armónica.
\end{proof}

\section{Circuitos}

De ahora en adelante vamos a distinguir dos conjuntos de nodos disjuntos que notaremos 
$A, Z$ y llamaremos \defn{entrada} y \defn{salida} respectivamente. Usualmente ambos serán solo un punto.

\begin{defi}[Voltaje]
  Dados $A, Z \subset \mathcal{X}$ disjuntos llamamos \defn{voltaje} a (la única) función de los vertices
  armónica en $(A \cup Z)^{c}$.
\end{defi}

\begin{obs}
  Es única dado que sabemos cuanto vale en $A, Z$.
\end{obs}

\begin{defi}[Corriente]
  Dado un voltaje $v$, definimos la $\defn{corriente}$ i
  como
  $$ i(x, y) = (v(x) - v(y))c(x, y).$$
\end{defi}

\begin{obs}
Notemos que la corriente es antisimétrica, es decir, $i(x, y) = -i(y, x)$
y, si $x \notin A\cup Z$
\begin{align*}
  \sum_{y \sim x} i(x, y) &= v(x) \sum_{y\sim x} c(x, y) + \sum_{y \sim x} v(y) c(x, y) \\
                          &= v(x) c(x) + c(x) \sum_{y \sim x} v(y) p(x, y) \\
                          &= 0. 
\end{align*}
Donde la ultima igualdad vale por ser $v$ armónica. Esto nos dice que el "flujo" neto 
en un punto es $0$, en analogía con el cálculo tradicional esta suma se suele notar 
como $\text{div }\theta(x).$
\end{obs}

\begin{defi}[Intesidad]
  Dado $a$ un punto de entrada o fuente, llamamos \defn{intensidad de corriente} 
  a la corriente que sale desde $a$, es decir,
  a la magnitud $\sum_{x \sim a} i(a, x)$  
\end{defi}

\begin{defi}[Flujo]
  Dados $A, Z$ conjuntos llamamos \defn{flujo} entre $A$ y $Z$ a una función $\theta : \mathcal{X} \times \mathcal{X} \to \R$ 
  tal que $\theta(x,y) = -\theta(y, x)$ para todo $x, y \in \mathcal{X}$ y 
  $\sum_{y \sim x} \theta(x , y) = 0$ para todo $x \notin A\cup Z$.
\end{defi}

A esta definiciones y las propiedades que vamos a mencionar a continuación se las 
llaman leyes en el contexto de la física.

\paragraph{Ley de Ohm} Si $x \sim y$ entonces 
$$i(x, y) r(x, y) = v(x) - v(y).$$

\paragraph{Ley de nodos de Kirchoff} La corriente $i$ es un flujo entre 
$A$ y $Z$.

\paragraph{Ley de ciclos de Kirchoff} Si $x_1 \sim x_2 \sim \dots \sim x_n = x_1$ entonces 
$$\sum_{k=0}^{n-1} i(x_k, x_{k+1})r(x_k, x_{k+1}) = 0.$$

\section{Probabilidad de escape y resistencia}

Volvamos a pensar en probabilidades, supongamos que tenemos $a, z \in \mathcal{X}$ un espacio
de estados con $a \neq 0$, pdoríamos preguntarnos cual es la probabilidad dado que comenzamos
en $a$ de llegar a $z$ antes de volver, esto es $\Prob_a[\tau_z < \tau_{a}^{+}]$. Comencemos 
pensando el problema partiendo de un estado $x$ cualquiera, vale que 
$$\E_x 1_{\{z\}}(X_{\tau_{\{a,z\}}}) = \Prob_x[\tau_z < \tau_{a}]$$
por el principio de existencia tenemos que esta es la única extensión armónica de 
$\ind{{\{a,z\}}}$. Cómo la función
$$ x \mapsto \frac{v(a) - v(x)}{v(a) - v(z)} $$
es armónica y vale lo mismo en los bordes tenemos que
$$ \Prob_x[\tau_z < \tau_{a}] = \frac{v(a) - v(x)}{v(a) - v(z)}.$$
Luego,
$$ 
\Prob_a[\tau_z < \tau_{a}^{+}] = \sum_{x \in V} P(a, x) \Prob_x[\tau_z < \tau_{a}] = 
\sum_{x \sim a} \frac{c(a, x)}{c(x)} \frac{v(a) - v(x)}{v(a) - v(z)} = \frac{\sum_{x \sim a} i(x, a)}{c(x) (v(a) - v(z))}.
$$
Donde la última igualdad vale por ley de Ohm. Este desarrollo que no parece tener un fin
claro motiva la siguiente definición.

\begin{defi}[Resistencia y conductancia efectiva]
  Definimos la \defn{resistencia efectiva} entre $a$ y $z$ como
  $$ \mathscr{R}(a \leftrightarrow z) = \frac{v(a) - v(z)}{\sum_{x \sim a} i(x, a)}.$$
  Además, definimos la \defn{conductancia efectiva} cómo su inverso multiplicativo y notamos
  $\mathscr{C}(a \leftrightarrow z).$
\end{defi}

\begin{obs}
  ¿De donde viene el nombre resistencia efectiva? Dada una red, si la reemplazamos
  por una con solo dos vertices, $a$ y $z$, unidos por una arista con resistencia 
  $\mathscr{R}(a \leftrightarrow z)$ si ahora aplicamos el mismo voltaje que a la original
  en estos dos puntos la intensidad de corriente va a ser la misma.
\end{obs}

\begin{defi}[Green]
  Dado un tiempo de parada $\tau$ definimos la \defn{función de Green} para un paseo
  al azar como
  $$G_\tau(a, x) = \E_a[\text{número de visitas antes de } \tau] = \E_a[\sum_{t\geq0} \ind{X_t = x, \tau > t}]$$
\end{defi}

\begin{obs}
  $G_z(a, a) = c(a) \mathscr{R}(a \leftrightarrow z)$ pues el número de visitas
  a $a$ partiendo desde $a$ antes de llegar a $z$ es una geométrica de 
  parámetro $\Prob_a[\tau_z < \tau_{a}^{+}].$ 
\end{obs}

\section{Simplificando redes}

\paragraph{Resistencias en serie} Dos aristas con resistencias $r_1, r_2$
en serie son equivalente a una única con resistencia $r_1 + r_2$. Más precisamente,
si $w \in V \setminus {A \cup Z}$ es un nodo de grado dos con vecinos $u_1, u_2$ y 
reemplazamos las aristas $(u_1, w), (w, u_2)$ por $(u_1, u_2)$ con resistencia
$r(u_1, w) + r(w, u_2)$, entonces los potenciales y corriente en $G\setminus\{w\}$
permanecen iguales y la corriente entre $u_1$ y $u_2$ es $i(u_1, w)$. 

\begin{figure}[ht]
  \centering
  \begin{tikzpicture}[
      line width=0.3pt,                      % trazo fino
      every node/.style={font=\small},
      pt/.style={circle,fill,inner sep=1.1pt}]
    % serruchito de resistencia: 3 picos y 3 valles (ancho total 1.2)
    \def\zz{-- ++(0.1,0.15) -- ++(0.2,-0.3) -- ++(0.2,0.3) -- ++(0.2,-0.3)
            -- ++(0.2,0.3) -- ++(0.2,-0.3) -- ++(0.1,0.15)}

    % --- Red original: u1 --- w --- u2 ---
    \node[pt,label=above:{$u_1$}] (u1) at (0,0) {};
    \node[pt,label=above:{$w$}]   (w)  at (2.82,0) {};
    \node[pt,label=above:{$u_2$}] (u2) at (5.64,0) {};
    \draw (u1) -- ++(0.81,0) \zz -- (w);
    \draw (w)  -- ++(0.81,0) \zz -- (u2);
    \node at (1.41,-0.5) {$r(u_1,w)$};
    \node at (4.23,-0.5) {$r(w,u_2)$};
    % el resto de la red sale de u1 y de u2
    \draw (u1) -- ++(-0.5,0.3);  \draw (u1) -- ++(-0.5,-0.3);
    \draw (u2) -- ++(0.5,0.3);   \draw (u2) -- ++(0.5,-0.3);
  \end{tikzpicture}
  \label{fig:serie}
\end{figure}

\begin{proof}
  La demostración consiste en ver que en la nueva red  $G \setminus \{w\}$ la corriente 
  y el voltaje siguen siendo igualess, esto lo hacemos chequeando la ley de Ohm y Kirchoff.
\end{proof}

\paragraph{Resistencias en paralelo} Dos aristas con conductancia $c_1, c_2$ que unen los mismos
vertices son equivalentes a una única con conductancia $c_1 + c_2$. Luego, 
todos los voltajes y corrientes se mantienen igual en $G \setminus \{e_1, e_2\}$
y la corriente $i(e) = i(e_1) + i(e_2)$. Esto vale para infinitas aristas en paralelo.

\begin{figure}[ht]
  \centering
  \begin{tikzpicture}[
      line width=0.3pt,                      % trazo fino
      every node/.style={font=\small},
      pt/.style={circle,fill,inner sep=1.1pt}]
    % serruchito de resistencia: 3 picos y 3 valles (ancho total 1.2)
    \def\zz{-- ++(0.1,0.15) -- ++(0.2,-0.3) -- ++(0.2,0.3) -- ++(0.2,-0.3)
            -- ++(0.2,0.3) -- ++(0.2,-0.3) -- ++(0.1,0.15)}

    % --- Red original: dos aristas en paralelo ---
    \node[pt,label={[xshift=-0.45cm]left:{$w_1$}}]  (w1) at (0,0) {};
    \node[pt,label={[xshift=0.45cm]right:{$w_2$}}] (w2) at (5,0) {};
    \draw (w1) to[out=90,in=180] (0.9,0.9) -- (1.9,0.9) \zz -- (4.1,0.9)
          to[out=0,in=90] (w2);
    \draw (w1) to[out=-90,in=180] (0.9,-0.9) -- (1.9,-0.9) \zz -- (4.1,-0.9)
          to[out=0,in=-90] (w2);
    \node at (2.5,1.3)  {$e_1$};
    \node at (2.5,-1.3) {$e_2$};
    % el resto de la red sale de w1 y de w2
    \draw (w1) -- ++(-0.5,0.3);  \draw (w1) -- ++(-0.5,-0.3);
    \draw (w2) -- ++(0.5,0.3);   \draw (w2) -- ++(0.5,-0.3);
  \end{tikzpicture}
  \label{fig:paralelo}
\end{figure}

\begin{proof}
  Chequear Ohm y Kirchoff.
\end{proof}

Veamos que efectivamente nos ayuda a resolver problemas con un ejemplo.

\begin{ejemplo}\label{ej:arbol-simetrico}
  Sea $\Gamma$ un árbol con raíz $\rho$ y sea $\Gamma_n$ su \defn{nivel $n$},
  es decir, el conjunto de vértices a distancia $n$ de $\rho$. Pedimos dos
  simetrías:
  \begin{enumerate}[(i)]
    \item todos los vértices de un mismo nivel tienen el mismo grado
          (\defn{árbol esféricamente simétrico});
    \item todas las aristas que unen el nivel $i-1$ con el nivel $i$ tienen la
          misma resistencia, que llamamos $r_i$, para $i \ge 1$.
  \end{enumerate}
  Por simetría, el voltaje es constante en cada nivel, así que pegar en un solo
  vértice todos los vértices de un nivel no altera las corrientes ni la
  resistencia efectiva. Después de pegar, entre el nivel $i-1$ y el nivel $i$
  quedan $\lvert\Gamma_i\rvert$ aristas en paralelo (una por cada vértice del
  nivel $i$), todas de resistencia $r_i$; en conjunto equivalen a una
  resistencia $r_i/\lvert\Gamma_i\rvert$. Estos tramos quedan en serie, y las
  resistencias en serie se suman:
  \begin{equation}\label{eq:reff-arbol}
    \mathscr{R}(\rho \leftrightarrow \Gamma_M)
      = \sum_{i=1}^{M} \frac{r_i}{\lvert\Gamma_i\rvert}.
  \end{equation}
  Como la conductancia total en $\rho$ es $c(\rho) = \lvert\Gamma_1\rvert / r_1$,
  la fórmula de la probabilidad de escape, $\Prob_\rho\{\tau_{\Gamma_M} < \tau_\rho^+\}
  = 1/\bigl(c(\rho)\,\mathscr{R}(\rho \leftrightarrow \Gamma_M)\bigr)$, da
  \begin{equation}\label{eq:escape-arbol}
    \Prob_\rho\{\tau_{\Gamma_M} < \tau_\rho^+\}
      = \frac{r_1/\lvert\Gamma_1\rvert}{\displaystyle\sum_{i=1}^{M} r_i/\lvert\Gamma_i\rvert}.
  \end{equation}
  Haciendo $M \to \infty$, el límite es positivo si y solo si la serie del
  denominador converge:
  \[
    \lim_{M\to\infty} \Prob_\rho\{\tau_{\Gamma_M} < \tau_\rho^+\} > 0
    \iff \sum_{i=1}^{\infty} \frac{r_i}{\lvert\Gamma_i\rvert} < \infty.
  \]
  En otras palabras, la caminata es transitoria exactamente cuando esa serie
  converge.
\end{ejemplo}

\section{Energía}

\begin{defi}[Energía]
  Definimos la \defn{energía} de un flujo $\theta$ como
  $$\mathcal{E}(\theta) = \sum_e \theta^{2}(e)r(e).$$
\end{defi}

\begin{teo}[Principio de Thomson]
  Dado un grafo finito y conexo,
  $$\mathcal{E}(\theta) = \inf \{ \mathcal{\theta} \colon \theta \text{ es un flujo unitario de } a \text{ a } z \}$$
  donde flujo unitario hace referencia a intensidad de corriente igual a uno.
\end{teo}

\section{Resultado principal} 

\begin{teo}[Criterio de Nash-Williams]\label{teo:principal}
  Si $\langle \Pi_n\rangle$ es una sucesión de cortes disjuntos que separan
  $a$ de $\infty$, entonces
  \[
    \mathscr{R}(a\leftrightarrow\infty) \;\ge\; \sum_n
    \Bigl(\sum_{e\in\Pi_n} c(e)\Bigr)^{-1}.
  \]
\end{teo}

\begin{proof}
  Por el principio de Thomson, $\mathscr{R}$ es el mínimo de la energía de
  un flujo unitario. Aplicando Cauchy--Schwarz en cada corte,
  \[
    1 \le \Bigl(\sum_{e\in\Pi_n}\abs{\theta(e)}\Bigr)^2
    \le \Bigl(\sum_{e\in\Pi_n} r(e)\,\theta(e)^2\Bigr)
        \Bigl(\sum_{e\in\Pi_n} c(e)\Bigr).
  \]
  Sumando sobre $n$ se obtiene la desigualdad.
\end{proof}

\begin{coro}
  Si la serie diverge, la caminata es recurrente.
\end{coro}

\begin{lem}\label{lem:aux}
  Sea $(X_n)$ una sucesión de variables $\iid$ con $\E\abs{X_1}<\infty$.
  Entonces $\frac1n\sum_{k\le n}X_k \asto \E[X_1]$.
\end{lem}

\begin{prop}
  Si $X_n \pto X$, entonces $X_n \dto X$.
\end{prop}

\begin{obs}
  Las flechas \verb|\dto|, \verb|\pto| y \verb|\asto| abrevian los tres
  modos de convergencia usuales.
\end{obs}

\begin{ejemplo}
  En $\Z^2$ los cortes concéntricos tienen $O(k)$ aristas, así que la serie
  armónica diverge y la caminata simple es recurrente.
\end{ejemplo}

\begin{ej}
  Probar que $\Z^3$ es transitoria construyendo un flujo de energía finita.
\end{ej}

\section{Comentarios finales}

Listas numeradas y con viñetas:
\begin{enumerate}[(a)]
  \item primer punto;
  \item segundo punto.
\end{enumerate}

\paragraph{Un párrafo con título} Texto que sigue en la misma línea.

\todo{completar esta sección}

%  Bibliografía
\begin{thebibliography}{99}
  \bibitem{LP} R.~Lyons y Y.~Peres, \emph{Probability on Trees and
    Networks}, Cambridge University Press, 2016.
  \bibitem{LPW} D.~A.~Levin, Y.~Peres y E.~L.~Wilmer, \emph{Markov Chains
    and Mixing Times}, 2.ª ed., American Mathematical Society, 2017.
\end{thebibliography}

\end{document}
